% ATTEMPT_STATUS: unresolved
% RESEARCH_GENERATED: 2026-10-01; GPT-6 Astra Ultra; individual mathematical attempt
% TOP_PROBLEM: {"id": 140, "title": "Equidistribution of Integer Multiples", "category_id": 1, "category": {"id": 1, "name": "number_theory", "display_name": "Number Theory", "description": "Properties of integers, prime numbers, Diophantine equations.", "slug": "number-theory", "order_index": 1, "created_at": "2026-07-31T15:26:25.670Z"}, "status": "open"}
% FIRST_POSTED: 2026-10-01
% Imported from: unsolvedmath_additional_500.tex; UnsolvedMath ID 140
% !TeX program = lualatex
% Compile twice: lualatex 140.tex
% Self-contained: no external bibliography, images, or \input files.
% Numeric section labels and visible numbers are original UnsolvedMath IDs.
\documentclass[11pt,a4paper]{article}
\usepackage[margin=24mm,headheight=14pt]{geometry}
\usepackage{fontspec}
\setmainfont{Latin Modern Roman}
\setsansfont{Latin Modern Sans}
\setmonofont{Latin Modern Mono}
\usepackage{amsmath,amssymb,mathtools}
\usepackage{unicode-math}
\setmathfont{Latin Modern Math}
\usepackage{microtype}
\usepackage{enumitem,xurl,needspace}
\usepackage[dvipsnames]{xcolor}
\usepackage{hyperref}
\hypersetup{unicode=true,colorlinks=true,linkcolor=MidnightBlue,urlcolor=MidnightBlue,
 pdftitle={UnsolvedMath Problem 140},pdfauthor={Source-annotated compilation},bookmarksnumbered=true}
\usepackage{bookmark,tocloft,titlesec,fancyhdr}
\setlength{\cftsecnumwidth}{4.2em}
\cftsetpnumwidth{2.5em}
\cftsetrmarg{4em}
\titleformat{\section}{\Large\bfseries\raggedright}{\thesection}{0.7em}{}
\pagestyle{fancy}\fancyhf{}
\fancyhead[L]{\small\sffamily UnsolvedMath research attempt}
\fancyhead[R]{\small Original catalogue IDs}
\fancyfoot[C]{\thepage}
\setlength{\parindent}{0pt}
\setlength{\parskip}{5pt plus 1pt}
\setlength{\emergencystretch}{3em}
\setcounter{tocdepth}{1}\setcounter{secnumdepth}{1}
\setlist[itemize]{leftmargin=3em,itemsep=4pt,topsep=4pt,parsep=0pt}
\allowdisplaybreaks
\newcommand{\N}{\mathbb{N}}
\newcommand{\Z}{\mathbb{Z}}
\newcommand{\Q}{\mathbb{Q}}
\newcommand{\R}{\mathbb{R}}
\newcommand{\C}{\mathbb{C}}
\newcommand{\F}{\mathbb{F}}
\newcommand{\A}{\mathbb{A}}
\newcommand{\PP}{\mathbb{P}}
\newcommand{\E}{\mathbb{E}}
\newcommand{\eps}{\varepsilon}
\newcommand{\dd}{\,\mathrm d}
\newcommand{\sref}[2]{\hyperlink{source:#1:#2}{[S#2]}}
\newif\ifshowcatalogue
\showcataloguetrue
\newcommand{\cataloguescope}[1]{\ifshowcatalogue
 \paragraph{Statement recorded in the supplied source.}
 #1\fi}
\begin{document}
% UnsolvedMath ID 140; source catalogue code GREEN-052

\setcounter{section}{139}

\section{Spreading Integer Multiples Around the Circle}\label{140}

{\small\sffamily UnsolvedMath ID 140 · Number Theory}

\subsection{Definitions and mathematical statement}

\paragraph{Shared notation.}Unless a section says otherwise, $\N=\{1,2,\ldots\}$,
$\Z$ is the ring of integers, $\Q,\R,\C$ are the rational, real, and complex
fields, $[N]=\{1,\ldots,\lfloor N\rfloor\}$, and $\log$ is the natural
logarithm. For real functions of a parameter tending to infinity,
$f\ll g$ and $f=O(g)$ mean $|f|\leq Cg$ eventually for some fixed $C$;
$f\gg g$ reverses this comparison; $f=o(g)$ means $f/g\to0$;
$f\sim g$ means $f/g\to1$; and $f\asymp g$ means both order bounds.
A subscript on $O$ or $\ll$ specifies allowed constant dependence.
The expression $n^{o(1)}$ in an upper bound means growth smaller than
$n^\epsilon$ for each fixed $\epsilon>0$. Local definitions govern any
exception, finite-field convention, or use of the same symbol for another
object. An asymptotic claim is not automatically an effective algorithm.



For $\theta\in\R$, count the multiset of images $\theta a\pmod1$, $a\in A$, on $\R/\Z$. Thus distinct integers producing the same image are still counted separately. Intervals are circular arcs of length $1/n$, with a consistent half-open endpoint convention. The parameter $c>0$ is fixed before the $n$-element integer set is chosen. \sref{140}{1} \sref{140}{2}

\paragraph{Statement recorded in the supplied source.}
Let $c > 0$ and let $A$ be a set of $n$ distinct integers. Does there exist $\theta$ such that no interval of length $\frac{1}{n}$ in $\mathbb{R}/\mathbb{Z}$ contains more than $n^c$ of the numbers $\theta a \pmod 1$, for $a \in A$? \sref{140}{1}

\paragraph{Residual task reported by the archive.}

The embedded review (2026-08-17) records: Lower the exponent threshold to arbitrary c>0 or find an obstruction. \sref{140}{1}

\subsection{Short English statement}

Can one dilation spread any finite set of integers around a circle so that no short arc receives more than a very small power of the set size?

\subsection{Research attempt}

\paragraph{Provenance and scope.}
This individual attempt was generated by GPT-6 Astra Ultra on 1 October 2026. The method was to derive explicit partial results and test the first proposed extension adversarially. The original statement and sources below are retained. No claim of mathematical novelty or complete solution is made.

\paragraph{Formal target and quantifiers.}
Let $A$ consist of $n$ distinct integers. For every fixed $c>0$, one seeks a dilation $\theta\in\R/\Z$ such that every arc of length $1/n$ contains at most $n^c$ of the points $\theta a$. Multiplicity is counted if different integers land at the same point. The primary problem list asks for arbitrarily small positive exponents. We derive a completely uniform second-moment bound, test it on arithmetic progressions, and identify why its collision budget cannot establish the requested subpower occupancy.

\paragraph{A random dilation controls the total close-pair count.}
Assume $n\geq3$, so $1/n<1/2$. Define
\[
 C(\theta)=\#\{(a,b)\in A^2:a\ne b,\ \|(a-b)\theta\|_{\R/\Z}\leq1/n\}.
\]
For every nonzero integer $d$, the map $\theta\mapsto d\theta\pmod1$ preserves uniform Lebesgue measure. Hence the probability of $\|d\theta\|\leq1/n$ is exactly $2/n$, with endpoints of measure zero. Linearity of expectation, requiring no independence between different differences, gives
\[
 \int_0^1C(\theta)\,d\theta=2(n-1).
\]
There exists $\theta$ with $C(\theta)\leq2(n-1)$. If an arc of length $1/n$ contains $m$ labelled dilates, every ordered pair among them has circular distance at most $1/n$. Thus
\[
 m(m-1)\leq C(\theta)\leq2(n-1),\qquad
 m\leq\frac{1+\sqrt{8n-7}}2.
\]
The same selected $\theta$ works simultaneously for all arcs; no union bound over an uncountable family is used. This proves an $O(\sqrt n)$ occupancy bound for every integer set, independently of the magnitude of its elements.

\paragraph{A finite bad-set estimate.}
For any integer $M\geq2$, if some arc contains at least $M$ points then $C(\theta)\geq M(M-1)$. Markov's inequality therefore yields
\[
 \operatorname{meas}\{\theta:\text{some }1/n\text{-arc contains }M\text{ points}\}
 \leq\frac{2(n-1)}{M(M-1)}.
\]
This is stronger than a bare existence statement when $M$ is a large multiple of $\sqrt n$: most dilations then work. However for $M=n^c$ with $c<1/2$, the bound becomes larger than one and has no content. Replacing Markov by the same calculation with a different constant cannot cure this exponent loss.

\paragraph{Structured sets show the estimate is far from a sharp obstruction.}
If $A=\{a_0,a_0+d,\ldots,a_0+(n-1)d\}$ with $d\ne0$, choose $\theta=1/(nd)$, interpreted modulo one. The resulting points are a translate of the equally spaced grid $\{j/n:0\leq j<n\}$. Every half-open arc of length $1/n$ contains exactly one point, and every closed arc contains at most two. Thus even sets with heavily repeated differences can admit an essentially optimal dilation. High additive structure is not itself a counterexample.

The endpoint convention does not affect the asymptotic question: for each fixed $c>0$, the constant two is at most $n^c$ for all sufficiently large $n$. The cases $n=1,2$ can be checked separately and are irrelevant to a large-$n$ exponent statement.

\paragraph{The missing higher-order information.}
A cluster of $M$ points spends order $M^2$ of the close-pair budget. The averaging argument only reduces that budget to order $n$, so it permits a cluster of order $\sqrt n$. To improve the exponent by moment methods, one would need estimates for simultaneous events involving many differences. Those events are correlated: for an arithmetic progression all differences are multiples of one step, and treating them as independent would be false. Conversely the progression example shows that such correlations can also be exploited constructively. A useful argument must distinguish arithmetic organisation rather than merely count pairs.

One precise possible input would be a bound on the measure of dilations creating a cluster of $n^c$ that is strictly less than one, uniformly over all distinct integer sets. The displayed Markov estimate is an explicit attempted bound for that quantity and explains exactly where it stops.

\paragraph{Verification and status.}
Distinctness of the integers is needed only to ensure every difference used in the measure calculation is nonzero. Negative integers and a zero element cause no problem. The common dilation is chosen after averaging the entire collision count, so all arcs are controlled at once. The established result is the uniform square-root bound and the associated bad-measure inequality. No exponent below one half is obtained by this method, and in particular the arbitrary-$c$ target remains unresolved.

\subsection{Sources}

{\small

\begin{itemize}

\item[\hypertarget{source:140:1}{S1}] ULAM.AI, \emph{UnsolvedMath}, supplied \texttt{problems.json}, record 140 (GREEN-052), statement and background. The embedded literature-review date is 2026-08-17; that is the archive’s review date, not a fresh certification. Dataset landing page: \url{https://huggingface.co/datasets/ulamai/UnsolvedMath}.

\item[\hypertarget{source:140:2}{S2}] B. Green, 100 Open Problems, updated January 2026, Problem 86 comments. \url{https://people.maths.ox.ac.uk/greenbj/papers/open-problems.pdf}. The author-hosted document was inspected on 27 September 2026; the archive’s local code is not a problem-number locator in this paper.

\end{itemize}

}

\label{end:140}
\end{document}
